\documentclass[aps,prd,reprint,nofootinbib,superscriptaddress]{revtex4-2}

\usepackage{graphicx}
\usepackage{hyperref}
\usepackage{amsmath, amssymb}
\usepackage{booktabs}
\usepackage{siunitx}

%% Bundle-skeleton header (Phase 7g flagship draft)
%% Bundle target: F (Tier 0, flagship review, ~80--150pp PRD-long-format)
%% Generated: 2026-05-04 (Phase 7c.7 close-out; substantive draft awaits
%%   the all-12-siblings-GREEN gate per Phase 7 Wave 13 sequencing).
%% Schema: docs/BUNDLE_DIRECTORY_SCHEMA.md
%% Source manifest: papers/F/source_manifest.md (41 source papers)
%% Synthesis brief: papers/F/SYNTHESIS_BRIEF.md
%% Change log: papers/F/change_log.md
%% Append log: papers/F/append_log.json (machine-readable)

% Canonical counts from scripts/update_counts.py.
% Auto-generated by update_counts.py — 2026-04-24T19:44:03.140400
% Do not edit manually. Run: uv run python scripts/update_counts.py
\newcommand{\totaltheorems}{3993}
\newcommand{\substantivetheorems}{3882}
\newcommand{\placeholdertheorems}{111}
\newcommand{\axiomcount}{1}
\newcommand{\sorrycount}{0}
\newcommand{\sorrypercent}{0.0}
\newcommand{\leanmodules}{167}
\newcommand{\leandefinitions}{3297}
\newcommand{\pythonmodules}{68}
\newcommand{\testfiles}{59}
\newcommand{\totaltests}{2836}
\newcommand{\figurecount}{111}
\newcommand{\notebookcount}{54}
\newcommand{\papercount}{19}
\newcommand{\aristotleproved}{322}
\newcommand{\aristotleruns}{44}
% --- Paper 18 / Phase 5t per-section counts (FermiHubbardDimer.lean) ---
\newcommand{\fhdTotal}{143}
\newcommand{\fhdWaveTwo}{13}
\newcommand{\fhdWaveThree}{6}
\newcommand{\fhdWaveFour}{22}
\newcommand{\fhdWaveFive}{22}
\newcommand{\fhdWaveSixABC}{38}
\newcommand{\fhdWaveSixExtra}{15}
\newcommand{\fhdWaveSeven}{19}
\newcommand{\fhdWaveEight}{8}


\begin{document}

\title{Fluid-Based Approaches to Fundamental Physics:\\
A Formally Verified Survey}

\author{John Roehm}

\date{\today}

\begin{abstract}
%% TODO: F abstract is the 13-bundle synthesis statement. Substantive
%% content lift awaits all-12-siblings-GREEN gate.
%%
%% Top-level synthesis claim NEW to F (per SYNTHESIS_BRIEF.md):
%%   The same antisymmetric-tetrad-bilinear (ADW) substrate that drives
%%   the emergent-gravity content of Bundle D3 (single Sakharov coefficient
%%   $G_N^{\rm Sak} = 12\pi/(N_f \Lambda_{\rm UV}^2)$ propagating across
%%   six observables) also drives the SK-EFT analog-Hawking content of
%%   Bundle D1 (across BEC, polariton, graphene; 92\% Lean-theorem reuse),
%%   the chirality-wall + modular-invariance content of Bundle D2
%%   (Spin$^{\mathbb{Z}_4}$ 5-bordism with $N_f = 16$ matching D3's
%%   Sakharov-coefficient $N_f$), the dark-sector NO-GO closures of
%%   Bundle D5 (causal-set + entropic + Jacobson-thermo-GR with full
%%   Phase 6m closure), and the topological quantum-computation
%%   foundations of Bundle D4 (one MTC, four faces). The substrate is
%%   one object; its twelve faces are thirteen bundles' worth of
%%   predictive register.

We consolidate twelve sibling-bundle threads of the SK--EFT Hawking
research program into a single, formally verified survey of the
substrate-condensation approach to fundamental physics. The
substrate's predictive register lives at the boundary partition
(\emph{cleared / falsified / structurally-blocked / reproduced}); this
review reports the substrate as a single object across all twelve
sibling threads. The synthesis claim, NEW to this flagship and absent
from any individual sibling, is that the same antisymmetric-tetrad-bilinear
(ADW) substrate that drives the emergent-gravity content of
Bundle~D3---a single Sakharov coefficient
$G_N^{\mathrm{Sak}} = 12\pi/(N_f\,\Lambda_{\mathrm{UV}}^{2})$
propagating across six observables---also drives the SK--EFT analog-Hawking
content of Bundle~D1 (across BEC, polariton, and graphene Dirac-fluid
platforms; $92\%$ Lean-theorem reuse), the chirality-wall and
modular-invariance content of Bundle~D2
($\mathrm{Spin}^{\mathbb{Z}_4}$ 5-bordism with $N_f = 16$ matching
D3's Sakharov-coefficient $N_f$), the dark-sector NO-GO closures of
Bundle~D5 (causal-set, entropic, and Jacobson thermodynamic-GR with
the full Phase~6m three-track closure), and the topological
quantum-computation foundations of Bundle~D4 (one MTC, four faces).
The substrate is one object; its twelve faces are thirteen bundles'
worth of predictive register. NO-GO results are reported as
first-class predictive content rather than absences: GW170817 falsifies
the vestigial-second-sound graviton at $\sim 7\times 10^{14}$; an
abelian-MTC falsifier closes the Bekenstein--Hawking entropy with the
Kaul--Majumdar $-3/2 \log A$ correction structurally rather than by
fiat; entropic-gravity dark energy is closed unanimously across $8/8$
candidates with quantitative Bayesian thresholds exceeding decisive.
The cosmological-constant problem is \emph{reproduced} (not solved) at
$\Lambda_{\mathrm{emerg}}/\Lambda_{\mathrm{obs}} \gtrsim 10^{100}$
(the rigorous Lean lower bound; heuristically $\sim 10^{122}$). All
load-bearing claims ship with Lean~4 modules at zero \texttt{sorry}
and one true axiom (\texttt{gapped\_interface\_axiom}; the 4+1D
gapped-interface conjecture).
\end{abstract}

\maketitle

\tableofcontents

%% =====================================================================
%% §1 — Strategic frame
%% =====================================================================
\section{Strategic frame}
\label{sec:strategic-frame}

\subsection{The substrate program in one paragraph}

The SK--EFT Hawking program asks a structural question: \emph{what
laboratory-substrate physics, when taken seriously as a fluid-EFT input,
reproduces or constrains Standard-Model-plus-General-Relativity
phenomenology at predictive precision?} The program's answer is the
ADW (Akama--Diakonov--Wetterich) tetrad-bilinear substrate~\cite{Akama1979,
Diakonov2011}: an
antisymmetric-tetrad-bilinear order parameter that condenses below a
microscopic scale $\Lambda_{\mathrm{UV}}$ and produces a Sakharov-style
emergent gravitational coupling
$G_N^{\mathrm{Sak}} = 12\pi/(N_f\,\Lambda_{\mathrm{UV}}^{2})$
through the heat-kernel $a_2$ coefficient of an
$N_f$-flavour Dirac fermion bundle~\cite{Sakharov1968,Adler1980,
Vassilevich2003}. Twelve sibling bundles in this series take that
substrate and read off its predictions across SK--EFT analog Hawking
radiation in three laboratory platforms (Bundle~D1, with companion
experimental letters E1 and E2), Standard-Model anomaly cancellation
and modular invariance (Bundle~D2 with PRL splash L2), emergent gravity
through black-hole thermodynamics (Bundle~D3 with PRL splashes L1 on
GW170817 and L3 on the BCH four laws), topological quantum-computation
foundations (Bundle~D4), the dark sector under substrate constraints
(Bundle~D5), and the methodology and verification stack itself
(Bundles~I1 and I2). This flagship review is the thirteenth bundle. It
synthesises the twelve sibling threads into a unified description of
the substrate's predictive register and reports the cross-thread
identities that none of the sibling bundles report individually.

\subsection{The 13-bundle publication architecture}

Per the project's strategy frame~\cite{Roehm2026Strategy}, the SK--EFT
Hawking program ships across $13$ publication targets organised by
target tier:

\begin{itemize}
\item \textbf{Tier 0 (1 paper):} this flagship review (Bundle~F).
\item \textbf{Tier 1 (5 deep papers):} Bundle~D1 (analog Hawking
  across three platforms; PRD long format), Bundle~D2 (anomaly
  constraints on Standard-Model particle content), Bundle~D3 (emergent
  gravity through BH thermodynamics; the heaviest single bundle at
  $\sim 50$pp), Bundle~D4 (topological quantum-computation foundations),
  Bundle~D5 (dark sector under substrate constraints).
\item \textbf{Tier 2 (3 PRL splashes):} Bundle~L1 (GW170817 falsifies
  the vestigial-second-sound graviton), Bundle~L2 (three SM fermion
  generations from modular invariance), Bundle~L3 (BCH four laws by
  regime).
\item \textbf{Tier 3 (2 infrastructure papers):} Bundle~I1 (verification
  methodology with worked cases), Bundle~I2 (verified statistics +
  Lean tensor-categories software paper).
\item \textbf{Tier 4 (2 experimental letters):} Bundle~E1 (Paris--LKB
  polariton companion to D1), Bundle~E2 (Dean--Kim--Lucas graphene
  companion to D1).
\end{itemize}

The architecture is intentionally append-friendly: new Tier-2 PRL
splashes can be added if a future Phase 6X wave produces a 4-page-PRL
shaped headline. Per~\cite{Roehm2026Strategy} \S 3.5, NO-GO results
are reported as first-class publishable content rather than absences.

\subsection{NO-GO results as first-class predictive content}

The substrate's predictive register lives at the \emph{boundary
partition} of phenomenology: each prediction lands in exactly one of
four categories.

\begin{enumerate}
\item \textbf{Cleared.} The substrate produces a quantitative
  prediction agreeing with observation within the published
  experimental band. Examples include the four cleared decision gates
  E.1--E.4 of the emergent-gravity content (D3 \S\ref{sec:emergent-gravity}),
  the polariton stimulated-Hawking-gain prediction $G > 0.5$ at
  $\omega/\kappa \approx 0.175$ (D1 with E1), and the strongest Phase~6m
  Track-C $f(R)$ candidates Exp and ArcTanh
  ($\Delta\!\mathrm{AIC} \simeq \Delta\!\mathrm{BIC} \gtrsim 20$,
  Plaza--Kraiselburd~\cite{Roehm2026D5}).
\item \textbf{Falsified.} The substrate produces a quantitative
  prediction in tension with observation at $\geq 5\sigma$ or
  Bayes-factor $|\log\mathcal{B}| \geq 5$ (Jeffreys-decisive). Examples
  include the GW170817 vestigial-graviton falsification at $\sim
  7\times 10^{14}$ (L1 with D3 \S\ref{sec:emergent-gravity}), the
  Wen--ADW factor-$6000$ deficit closing the
  emergent-non-Abelian-gauge-symmetry pathway (D3 \S\ref{sec:gauge-erasure}),
  the abelian-MTC falsifier on Bekenstein--Hawking entropy with the
  Kaul--Majumdar $-3/2\log A$ correction (D3 \S\ref{sec:emergent-gravity}
  with D4 \S\ref{sec:topological-qc}), and the unanimous $8/8$
  Phase~6m Track-B entropic-gravity dark-energy closure with all
  Bayesian thresholds exceeding decisive (D5 \S\ref{sec:dark-sector}).
\item \textbf{Structurally blocked.} The substrate's mechanism is
  inapplicable for a load-bearing structural reason that is independent
  of the observational target. Examples include the Gibbs--Duhem
  obstruction theorem~\cite{Roehm2026D5} ruling out single-scalar
  self-tuning emergent-vacuum frameworks at the kinematic level, the
  fracton diffeomorphism-invariance no-go (D3
  \S\ref{sec:gauge-erasure}, citing~\cite{PretkoRadzihovsky2018}),
  and the four-factor orthogonality principle on the dark-energy
  obstruction (D5 \S\ref{sec:dark-sector}).
\item \textbf{Reproduced (not solved).} The substrate produces a
  prediction that matches observation at the structural level (sign,
  scale, or qualitative behaviour) but fails to solve the corresponding
  fine-tuning problem. The leading example is the cosmological-constant
  problem: the Sakharov coefficient predicts an effective $\Lambda$
  that overshoots the observed value by
  $\Lambda_{\mathrm{emerg}}/\Lambda_{\mathrm{obs}} \gtrsim 10^{100}$
  (rigorous Lean lower bound) or heuristically $\sim 10^{122}$. The
  substrate inherits the $a_0$ coefficient's universality across
  emergent-gravity proposals at one-loop; emergent gravity does not
  shelter from the cosmological-constant problem.
\end{enumerate}

The substrate's predictive register is the union of these four
partitions. \S\ref{sec:substrate-register} is the consolidated table.

\subsection{Substrate-identity synthesis claim NEW to this flagship}

Each sibling bundle reports its own thread; this flagship reports the
substrate as one object. The synthesis claim NEW to F---absent from any
individual sibling bundle and load-bearing for the architectural
argument---is that all twelve sibling threads share the same ADW
substrate. Concretely:

\begin{itemize}
\item The single Sakharov coefficient
  $G_N^{\mathrm{Sak}} = 12\pi/(N_f\,\Lambda_{\mathrm{UV}}^{2})$
  derived from the heat-kernel $a_2$ coefficient on the substrate's
  $N_f$-flavour Dirac fermion bundle (D3 \S\ref{sec:emergent-gravity})
  uses the same $N_f = 16$ count that, on the modular-invariance side,
  uniquely fixes the Standard-Model anomaly classification via
  $\mathrm{Spin}^{\mathbb{Z}_4}$ 5-bordism (D2
  \S\ref{sec:chirality-wall-modular} with L2). The $N_f$ in the
  Sakharov coefficient and the $N_f$ in the modular-invariance
  derivation of three SM fermion generations are the same $N_f$.
\item The same SK--EFT axioms, KMS condition, and
  fluctuation-dissipation relation that close the analog-Hawking
  spectrum on the BEC platform (D1 \S\ref{sec:analog-hawking})
  re-close on the polariton (E1) and graphene Dirac-fluid (E2)
  platforms with $\sim 92\%$ Lean-theorem reuse. The substrate is
  platform-agnostic at the formal level; the platforms supply only
  the dispersion data.
\item The MTC that carries the SU(2)$_k$ Chern--Simons logarithmic
  $-3/2\log A$ correction to BH entropy (D3
  \S\ref{sec:emergent-gravity}) is the same MTC that supports
  Hayden--Preskill information recovery on the horizon (D4
  \S\ref{sec:topological-qc}) and that underlies the Drinfeld-centre
  cross-bridge to the Standard-Model anomaly classification (D2
  \S\ref{sec:chirality-wall-modular} with D4
  \S\ref{sec:topological-qc}). One MTC, four faces (Drinfeld centre +
  WRT partition + HP-QEC + RT/CH knife-edge), per D4's synthesis
  claim.
\item The dark-sector NO-GO closures of D5
  (\S\ref{sec:dark-sector})---across causal-set, entropic-gravity, and
  Jacobson thermodynamic-GR mechanism families---use the \emph{same}
  Sakharov-induced-gravity criterion that distinguishes substrate
  classes within Phase~6m's Track C. The Sakharov four-condition
  criterion validated on Volovik--Jannes ${}^3\mathrm{He}$-A and
  falsified on Finazzi--Liberati--Sindoni acoustic-BEC is, to our
  knowledge, the first systematic $\Lambda_J$-vs-$\Lambda_{HK}$
  comparison on a common substrate in the literature.
\end{itemize}

The substrate is one object. Its twelve faces are thirteen bundles'
worth of predictive register. The cross-thread identities listed above
are the F-unique synthesis content, structurally absent from any
individual sibling bundle.

%% =====================================================================
%% §2 — Architectural scope
%% =====================================================================
\section{Architectural scope}
\label{sec:architectural-scope}

This section lifts the architectural-scope statement from
\texttt{docs/ARCHITECTURE\_SCOPE.md} (the canonical project-scope
document). The text is reproduced essentially verbatim with light
typographic adjustment; the source document supersedes this section
on any disagreement and is authoritative as of the flagship-submission
date~\cite{Roehm2026Strategy}.

\subsection{The three layers}

The SK--EFT Hawking program is organised around a three-layer
emergent-physics architecture:

\begin{itemize}
\item \textbf{Layer 1.} Discrete lattice Hamiltonian plus a Dirac
  fermion bundle. Canonical Lean modules:
  \texttt{LatticeHamiltonian}, \texttt{FermionBag4D}.
\item \textbf{Layer 2.} Coarse-grained effective field theory in the
  Schwinger--Keldysh closed-time-path formalism. Canonical modules:
  \texttt{AcousticMetric}, \texttt{SKDoubling},
  \texttt{SecondOrderSK}.
\item \textbf{Layer 3.} Emergent gauge fields plus emergent gravity
  produced by the ADW (Akama--Diakonov--Wetterich) tetrad-bilinear
  condensation mechanism, augmented by gauge erasure, the chirality
  wall, and vestigial-gravity intermediate ordering. Canonical modules:
  \texttt{ADWMechanism}, \texttt{GaugeErasure}, \texttt{ChiralityWall},
  \texttt{VestigialGravity}.
\end{itemize}

\subsection{Predictive scope by layer}

\paragraph{Layer 3, SM and GR sector (IN SCOPE).} The architecture
commits to Standard-Model-plus-General-Relativity emergent physics
under the tested mechanisms: ADW tetrad condensation producing
emergent gravity, gauge erasure producing $U(1)$ from a non-Abelian
lattice, fracton hydrodynamics producing subdiffusive emergent
currents, and vestigial gravity as a distinct intermediate ordering.
Layer-3 predictions in this sector are falsifiable, formalised
(currently $\sim \totalsubstantivetheorems$ Lean theorems with zero
\texttt{sorry} and one true axiom; see \S\ref{sec:methodology}), and
actively tested against laboratory platforms (BEC, polariton, graphene
Dirac fluid; D1, E1, E2).

\paragraph{Dark-sector physics (out-of-scope under tested
mechanisms).} The Phase 5y closure verdict (2026-04-23) reports that
after six deep-research rounds of Volovik-family dark-energy
mechanisms (q-theory in four realisations, vestigial-gravity equation
of state, second-sound graviton), all returned NO-GO for DESI~DR2
compatibility. The structural obstruction is formalised in three Lean
modules: \texttt{GibbsDuhemTheorem.lean} (any single-scalar
self-tuning emergent-vacuum framework with Gibbs--Duhem equilibrium
locks $w_{\mathrm{vac}} = -1$), \texttt{QTheoryNoGoTheorem.lean} (the
obstruction is realisation-independent across the four tested KV
q-theory constructions), and
\texttt{DarkEnergyObstructionPrinciple.lean} (four-factor orthogonality
diagnoses why vestigial gravity also fails:
$\text{GD} \cap c_s^{2}\!\geq\!0 \cap \text{natural }T_c \cap
\text{MICROSCOPE}$).

The Phase 5y closure is a sharpening of scope, not a retreat: the
architecture remains intact for SM-plus-GR physics; it is explicit
about not predicting cosmological-constant-type physics under
mechanisms currently tested. The Phase 6m closure (2026-04-30; see
below) extends the scope statement to three further mechanism families.

\paragraph{What this does not claim.} The Phase 5y closure does not
claim that q-theory is falsified as a vacuum-stability framework
(only as a late-time dark-energy framework); does not claim that
vestigial gravity is ruled out as a concept (the H4 framework
PARTIAL verdict remains; the ordering and $\mathbb{Z}_4$ structure
are mathematically consistent and may find applications elsewhere);
does not claim that no emergent dark energy is possible (the
obstruction covers Volovik-family mechanisms specifically; Phase 6m
extends to entropic-gravity and Jacobson-thermo-GR mechanism families
below); and does not falsify the three-layer architecture (only the
dark-sector predictive claim is scoped out).

\subsection{Phase 6m closure (2026-04-30)}

After three independent six-round deep-research probes of the mechanism
families flagged out-of-scope by Phase 5y (Tracks A, B, C in Phase~6m),
Phase~6m closes at the Lean-formalisation scope. Project totals:
$\sim \totalsubstantivetheorems$ substantive theorems across
\totalmodules{} modules, $\axioms$ axiom, $\sorries$ \texttt{sorry}.

\textbf{Track-level verdicts.}
\begin{itemize}
\item \textbf{Track A --- Causal-set DE} (\texttt{CausalSetDarkEnergy.lean},
  15 theorems): three NO-GO-R5 phenomenological + causal-set
  d'Alembertian NO-GO-R2 reaffirmed via 4D gradient instability.
  Three publishable structural caveats survive independent of DESI:
  GD-inapplicability robust under all four sprinkling prescriptions;
  Barrow-bound prescription dependence; BDG
  $\sigma_\Lambda = \alpha_{\mathrm{BDG}}/\sqrt{V}$ first-principles
  decomposition.
\item \textbf{Track B --- Entropic-gravity DE}
  (\texttt{EntropicGravityDarkEnergy.lean}, 14 theorems): \emph{eight
  NO-GO-R5 unanimous} --- the first complete-mechanism-family NO-GO
  closure in Phase~6m. Verlinde~2017, Padmanabhan-CosMIn,
  Hossenfelder--Verlinde post-Yoon--Guha, Cadoni--Tuveri DEC, Li~2004
  HDE, Tsallis HDE, Barrow HDE, and Odintsov-D'Onofrio-Paul
  $\Omega_k$ all closed with quantitative Bayesian thresholds
  exceeding Jeffreys-decisive ($|\!\log\mathcal{B}| \geq 5$). Seven of
  eight candidates have $r_d$-independent NO-GO mechanisms.
\item \textbf{Track C --- Jacobson thermodynamic-GR DE}
  (\texttt{JacobsonThermoGRDarkEnergy.lean}, 12 theorems):
  highest-survival track with five-plus R5 survivors. M3 EGJ $f(R)$
  Exp and ArcTanh are the strongest CLEARED-R5 of any track in
  Phase~6m (Plaza--Kraiselburd:
  $\Delta\!\mathrm{AIC} \simeq \Delta\!\mathrm{BIC} \gtrsim 20$,
  Jeffreys ``very strong''); Hu--Sawicki NO-GO-R5 via the chameleon
  Solar-System constraint; Pure Lovelock NO-GO-R5 reaffirmed at the
  $1\sigma$-box edge of Quintom-B; M8 KSS CLEARED-R5 conditional via
  the path-(a) Arata--Liberati--Neri reduction.
\end{itemize}

\textbf{Phase 6e cross-bridge final closure.} The Sakharov
four-condition criterion biconditional with $\Lambda_J = \Lambda_{HK}$
is validated on Volovik--Jannes ${}^3\mathrm{He}$-A and falsified on
Finazzi--Liberati--Sindoni acoustic-BEC (condition (ii) fails---only
phonons have a BEC effective metric). To our knowledge this is the
first systematic $\Lambda_J$-vs-$\Lambda_{HK}$ comparison on a common
substrate in the literature.

\textbf{Wave 4 unified Phase~6m 7-class GD taxonomy}
(\texttt{DarkSectorClassificationExtension.lean}, 10 theorems):
consolidates the three-track verdicts into a single classification
namespace with a 3-tier GD-applicability gradient. M8 KSS uniquely
populates Class~(a). The unimodular reformulation as
$\Lambda_{HK}$-escape route admits 5-of-6 Track-C survivors except
KSS.

\subsection{Predictive-scope recalibration after Phase 6m}

The Layer-3 dark-energy scope statement, sharpened from the Phase 5y
formulation:
\begin{itemize}
\item \textbf{NO-GO (machine-checked):} Volovik-family mechanisms;
  Track A causal-set DE phenomenological; Track B entropic-gravity DE
  complete-class; Track C M3 Hu--Sawicki + M4 Pure Lovelock; Track A
  causal-set d'Alembertian.
\item \textbf{PARTIAL-VIABLE (R5 cleared, requires monitoring):}
  Track C M1 Jacobson 1995, M2/M7 Padmanabhan/CosMIn (epistemic flag),
  M9 Volovik--Jannes (under fixed-$r_d$).
\item \textbf{CLEARED-R5 strongest (preferred over $\Lambda$CDM by
  $\Delta\!\mathrm{AIC} \gtrsim 20$):} Track C M3 EGJ $f(R)$
  Exponential and ArcTanh; Starobinsky marginal.
\item \textbf{CONDITIONAL CLEARED-R5 / OPEN-R6+:} Track C M8 KSS via
  path-(a) PARTIAL closure; OPEN via path-(c) universal-horizon
  thermodynamics.
\item \textbf{Three publishable structural caveats} (Track A,
  independent of DESI): GD-inapplicability across prescriptions;
  Barrow-bound prescription dependence; BDG $\sigma_\Lambda$
  first-principles decomposition.
\end{itemize}

\subsection{Architectural implications}

\begin{enumerate}
\item \textbf{Program value preserved.} The Lean backbone of the
  program continues to build cleanly across all 13 sub-phases.
  $\sim \totalsubstantivetheorems$ theorems, $\sorries$ \texttt{sorry},
  $\axioms$ axiom.
\item \textbf{Scope explicit.} The predictive-scope boundary is
  formally stated in Lean
  (\texttt{ClassificationTableDark.lean},
  \texttt{DarkSectorClassificationExtension.lean}) and consolidated
  here. Future Phase 6X / Phase 7 architectural conversations can
  reference this scope explicitly without ambiguity.
\item \textbf{GO/NO-GO methodology validated.} Phase 5y demonstrated
  that the program can execute a six-round deep-research probe, land
  a structural NO-GO, and harvest the result as formal Lean content
  without reopening or re-scoping the underlying decision. Phase 6m
  validated the methodology on three further mechanism families. This
  is a reusable methodology for future high-leverage architectural
  bets.
\item \textbf{NO-GO-as-publishable-content.} The three-track Phase 6m
  closure ships as load-bearing content of Bundle~D5
  \S\ref{sec:dark-sector}, not as background scope-disclaimer. The
  publishable-novelty claim of Bundle~D5 is the first
  complete-mechanism-family NO-GO closure in Phase 6m for the
  entropic-gravity dark-energy class.
\end{enumerate}

%% =====================================================================
%% §3 — Gauge erasure (paper3)
%% =====================================================================
\section{Gauge erasure}
\label{sec:gauge-erasure}

A higher-form-symmetry analysis of the Glorioso--Liu Schwinger--Keldysh
effective-field-theory programme~\cite{GloriosoLiu2018} establishes
that the hydrodynamic layer \emph{erases} non-Abelian one-form gauge
symmetry: only $U(1)$ survives the fluid-cell coarse-graining as a
continuous one-form. Non-Abelian gauge fields survive only as discrete
centre symmetries that re-enter through the deconfinement transition
(Bundle~D3 \S 14, with cross-bridge to D4 \S 4 below). The erasure
theorem (\texttt{GaugeErasure.gauge\_erasure},
\texttt{GaugeErasure.u1\_survival}) is the first of three
predictive-boundary cuts that force the ADW substrate route to the
tetrad-condensation channel; the other two are the perturbative
Wen-ADW deficit and the GW170817 vestigial-graviton falsification
(\S\ref{sec:emergent-gravity}).

\subsection{Erasure mechanism}

Non-Abelian one-form symmetries require a colour-charged operator
algebra on each fluid cell. Hydrodynamic coarse-graining washes out
the colour correlation outside the screening length while preserving
the Abelian charge, because the Abelian charge is additive and
survives averaging. The effective fluid Lagrangian therefore admits
only $U(1)$ continuous one-forms; non-Abelian remnants survive only
as the centre $\mathbb{Z}_N \subset \mathrm{SU}(N)$, which is
gauge-trivial in the fluid limit but reappears through the
Polyakov-loop modulus in the deconfinement order parameter
(see~\cite{Polyakov1978,SvetitskyYaffe1982}).

The erasure theorem is the substrate's structural answer to a
question that frames much of fluid-EFT model-building: \emph{why does
a fluid layer not naturally produce non-Abelian gauge fields?} The
answer is that, at the formal level, it cannot: the higher-form
selection rule is a hydrodynamic coarse-graining identity, not an
empirical regularity.

\subsection{Universality across substrates}

The erasure mechanism is platform-universal in a sense that this
flagship's substrate-identity claim leans on. Relativistic
hydrodynamics, Pretko--Radzihovsky fracton
hydrodynamics~\cite{PretkoRadzihovsky2018}, and Maldacena's
$\mathcal{N}=4$ super-Yang--Mills hydrodynamic
limit~\cite{Maldacena1997} all share the same erasure selection
rule. Domain walls---not Goldstone bosons---carry the residual
non-Abelian information. This links forward to the Polyakov-loop
confinement order parameter in Bundle~D3 \S 14 and to the colour-flavour
locked diquark condensate in D3 \S 16 with~\cite{AlfordRajagopalWilczek1999}.

The platform-universality of erasure is what turns the substrate's
SK--EFT analog Hawking content of D1 (BEC, polariton, graphene Dirac
fluid) into a single substrate-level claim rather than three
platform-specific claims: each platform's hydrodynamic fluid is
governed by the same erasure rule, so the SK--EFT axioms close on
each platform with the same continuous-$U(1)$ residual.

\subsection{Decisive test: $\mathcal{N}=4$ SYM}

$\mathcal{N}=4$ super-Yang--Mills carries a one-form $\mathrm{SU}(N)$
symmetry above the conformal threshold but a $\mathbb{Z}_N$ centre at
the hydrodynamic fixed point. The erasure theorem predicts the centre
survives and the continuous $\mathrm{SU}(N)$ one-form does not; this
is the substrate's decisive test against an alternative
continuous-non-Abelian hydrodynamic regime. The prediction is
independent of holographic-duality assumptions: the
hydrodynamic-fixed-point limit of any candidate gauge theory selects
$U(1)$-plus-centre at the SK-effective level.

\subsection{Cross-bridge to D4 \S 4 (Drinfeld centre)}

The discrete-centre $\mathbb{Z}_N$ residual from gauge erasure is
not a phenomenological remnant: it carries the substrate's
topological-quantum-computation data. Specifically, the $\mathbb{Z}_N$
centre of $\mathrm{SU}(N)$ enters the modular tensor category of
Bundle~D4 \S 4 through the Drinfeld-centre construction
$\mathcal{Z}(\mathrm{Vec}_{\mathbb{Z}_N})$, which is the canonical
$3+1$D topological order anchoring the Standard-Model anomaly
classification of Bundle~D2 (\S\ref{sec:chirality-wall-modular}).
The cross-bridge \emph{D3~\S 3 $\to$ D4~\S 4 $\to$ D2~\S 3.4} is one
substrate-identity thread: erasure selects the centre, the centre
populates the Drinfeld category, and the Drinfeld category supplies
the SM anomaly cancellation. The substrate-identity synthesis claim
is the observation that this thread closes self-consistently across
three sibling bundles without an additional mechanism.

%% =====================================================================
%% §4 — SK-EFT analog Hawking radiation across three platforms (D1 + E1 + E2)
%% =====================================================================
\section{SK--EFT analog Hawking radiation across three platforms}
\label{sec:analog-hawking}

%% TODO: 12pp lift from Bundle D1 (full content) + companion summaries from E1, E2.
%% Key claims:
%%   - Three SK-EFT axioms + KMS + FDR closure in SKDoubling.firstOrder_KMS_optimal.
%%   - Three platforms (BEC, polariton, graphene) with 92\% Lean theorem reuse.
%%   - Exact WKB connection formula with three non-perturbative effects.
%%   - Polariton stimulated-Hawking gain G > 0.5 at omega/kappa = ln(3)/(2\pi) approx 0.175 with N_probe = 100 photons.
%%   - Graphene T_H approx 2.4 K, delta_disp approx -2.8\%, Wiedemann-Franz violation L/L_0 > 200.

\textit{Substantive content awaiting all-12-siblings-GREEN gate (D1 + E1 + E2 already closed).}

%% =====================================================================
%% §5 — Chirality wall + modular invariance (D2 + L2)
%% =====================================================================
\section{Chirality wall and modular invariance}
\label{sec:chirality-wall-modular}

Bundle~D2~\cite{Roehm2026D2} consolidates the substrate's
chirality-wall + modular-invariance content into a single deep paper
treating the Standard-Model fermion content as a $\mathrm{Spin}^{\mathbb{Z}_4}$
5-bordism anomaly classification with $\mathbb{Z}_{16}$ periodicity.
The PRL splash Bundle~L2~\cite{Roehm2026Modular} extracts the
headline four-page result on the modular-invariance derivation of
three Standard-Model fermion generations.

\subsection{Modular framing and the central charge}

The substrate's lattice fermion content carries a central charge
$c$ that, on the modular-framing constraint, is required to be a
consistent boundary condition for a $(2{+}1)$D topological field
theory. The modular generation constraint (D2 \S 2) shows that with
the Standard-Model fermion content (15 Weyl fermions per generation,
plus the right-handed neutrino $\nu_R$) the central charge is
$c = 16 g$ where $g$ is the number of generations. The 16-convergence
follows from the Rokhlin theorem on $\mathrm{Spin}^{\mathbb{Z}_4}$
4-manifolds: the signature divides by $16$, forcing $N_f = 16$ per
generation as the unique consistent count.

Without $\nu_R$, the central charge produces a fractional anomaly that
fails to lift to a consistent topological boundary condition. With
$\nu_R$, the lift is unique up to discrete data captured by the
$\mathbb{Z}_{16}$ anomaly classification.

\subsection{The Standard-Model $\mathbb{Z}_{16}$ anomaly}

The $\mathbb{Z}_{16}$ generator of the $\mathrm{Spin}^{\mathbb{Z}_4}$
5-bordism group acts on the substrate's per-generation fermion
content as the unique non-trivial generator. A single Standard-Model
generation contributes a $\mathbb{Z}_{16}$ charge of $1$
(\texttt{dai\_freed\_spin\_z4}); three generations stack to a
$\mathbb{Z}_{16}$ charge of $3$, which is anomaly-canceling modulo
the SPT-stacking 16-fold periodicity. The cross-section consistency
of the three generations against the $\mathbb{Z}_{16}$ anomaly is the
substrate's structural answer to the Standard-Model fermion-generation
counting problem: \emph{three generations is the unique fermion-content
solution that anomaly-cancels modulo the $\mathbb{Z}_{16}$ classification.}

The minimal Drinfeld-centre anchor (D2 \S 3.4) shows that the
$\mathbb{Z}_{16}$ anomaly classification arises from a Drinfeld
centre $\mathcal{Z}(\mathrm{Vec}_{\mathbb{Z}_4})$ on the substrate's
toric-code-style boundary condition. This is the cross-bridge to
Bundle~D4 \S 4, where the same Drinfeld-centre construction supplies
the substrate's topological-quantum-computation infrastructure
(\S\ref{sec:topological-qc}). The Drinfeld-centre object on the D4
side and the Drinfeld-centre object on the D2 side are the same
algebraic object in the substrate's category-theoretic data.

\subsection{The chirality wall: three-pillar synthesis}

The chirality-wall construction synthesises three independent pillars:

\paragraph{Pillar 1 (Golterman--Shamir no-go and TPF evasion).} A
chiral lattice gauge theory with manifest gauge invariance cannot
exist by the Golterman--Shamir no-go; the substrate's chirality wall
evades this obstruction through a topological-phase-factor (TPF)
construction that breaks gauge invariance manifestly while preserving
the topological invariant.

\paragraph{Pillar 2 (Gioia--Thorngren positive construction).} A
positive construction of the chirality wall is provided by the
Gioia--Thorngren mechanism: a 4+1D gapped interface between a trivial
$\mathbb{Z}_2$-symmetry-protected topological (SPT) phase and a
non-trivial $\mathbb{Z}_4$-anomalous boundary supports the
Standard-Model chiral fermion content as the boundary excitation
spectrum.

\paragraph{Pillar 3 (Onsager + Witten anomaly).} The Onsager-style
algebraic-geometry construction recovers the Witten anomaly as the
unique global obstruction to the Standard-Model fermion content
without $\nu_R$. With $\nu_R$ included, the anomaly cancels at the
classification level but persists as a tracked physical signature in
the $\theta$-angle data of Bundle~D3 \S 13 (Phase~6l Z16 strong-CP
substrate-silence verdict).

The three pillars synthesise into the chirality-wall theorem: the
4+1D gapped-interface axiom (\texttt{gapped\_interface\_axiom}) is
the unique structural ingredient required, and the substrate's
fermion content is the unique 4-manifold-boundary spectrum consistent
with the axiom.

\subsection{The single axiom \texttt{gapped\_interface\_axiom}}

The substrate's Bundle~D2 ships with one true axiom:
\texttt{gapped\_interface\_axiom}, the 4+1D conjecture that asserts
the existence of the gapped interface required by Pillar 2. This is
the only axiom in the entire 13-bundle program that has not been
discharged at the Lean-formalisation scope. All other load-bearing
substrate claims are theorems with zero \texttt{sorry}.

The axiom's status is partial in a specific sense: it is a
mathematical conjecture (existence of a particular gapped-interface
construction) that mathematicians have not yet proved, rather than a
physical postulate. Its discharge is a Mathlib-upstream cycle target
(see Bundle~I2~\cite{Roehm2026Strategy}) and is on the Phase 8 outlook
roadmap. Bundle~D4 inherits the axiom but introduces no new ones
(\S\ref{sec:topological-qc}).

\subsection{The L2 PRL splash}

Bundle~L2~\cite{Roehm2026Modular} extracts the four-page PRL splash
from the modular-invariance pillar (D2 \S 2) and ships as a stand-alone
result on \emph{three Standard-Model fermion generations from modular
invariance}. The splash carries the central-charge-modulo-16 derivation,
the Rokhlin 16-convergence, and the cross-section consistency check;
it does not carry the full $\mathrm{Spin}^{\mathbb{Z}_4}$ 5-bordism
anomaly classification (which lives in D2). The L2 splash and the
D2 deep paper share the same Lean theorem anchor
(\texttt{Z16Classification.dai\_freed\_spin\_z4}) character-for-character.

\subsection{Cross-bridge to D3's $N_f$}

The substrate-identity synthesis claim of this flagship leans on the
observation that the $N_f = 16$ derived in D2 (modular framing /
Rokhlin / 16-convergence) and the $N_f$ in the Sakharov coefficient
$G_N^{\mathrm{Sak}} = 12\pi/(N_f\,\Lambda_{\mathrm{UV}}^{2})$ derived
in D3 \S 17 (heat-kernel calibration on the substrate's Dirac-fermion
bundle) are the same $N_f$. The same fermion-content count that fixes
the Standard-Model anomaly classification also fixes the substrate's
emergent gravitational coupling. This is one of the four substrate-identity
threads listed in \S\ref{sec:strategic-frame} and is the F-unique
synthesis content that none of the sibling bundles report individually.

%% =====================================================================
%% §6 — Emergent gravity through BH thermodynamics (D3 + L1 + L3)
%% =====================================================================
\section{Emergent gravity through black-hole thermodynamics}
\label{sec:emergent-gravity}

%% TODO: 30pp lift from Bundle D3 (full content) + L1 splash + L3 splash.
%% Heaviest single F section; main substrate-identity content.
%% Key claims:
%%   - Single-coefficient identity G_N^Sak = 12pi/(N_f Lambda_UV^2) propagating across 6 observables.
%%   - Three closures forcing ADW: Wen-deficit, fracton no-go, GW170817 falsification.
%%   - Four cleared decision gates: E.1, E.2, E.3, E.4.
%%   - BCH four laws partitioned by regime mass M_c.
%%   - BH entropy with Kaul-Majumdar -3/2 log A; Sen 4D Schwarzschild -45/2 disagreement; abelian-MTC falsifier.
%%   - CC reproduced (not solved) at >10^100 (heuristic ~10^122).
%%   - EC torsion 2.05e-77 GeV ~46 orders below KRT.
%%   - 5 algebraic-Lorentzian-geometry first-formalisation appendices.

\textit{Substantive content awaiting D3 Stage-13 GREEN gate (in flight at draft creation; verdict pending).}

%% =====================================================================
%% §7 — Topological QC foundations (D4)
%% =====================================================================
\section{Topological quantum-computation foundations}
\label{sec:topological-qc}

Bundle~D4 (\cite{Roehm2026D4}) consolidates the substrate's
topological-quantum-computation infrastructure into a single
machine-verified account. To our knowledge, D4 ships the first
machine-checked construction of the Drinfeld quantum groups
$U_q(\mathfrak{sl}_2)$ and $U_q(\mathfrak{sl}_3)$ in any proof
assistant, the first complete machine-checked Ising modular-tensor-category
(MTC) braided structure with all six hexagon equations and ribbon
equations, the Witten--Reshetikhin--Turaev partition function on
$S^{2}\times S^{1}$ and surgery formula on three-manifolds, the
doublon-SWAP gate as the substrate's first symmetry-protected
topological quantum gate, the Hayden--Preskill error-correcting-code
structure on horizon MTCs, and the Ryu--Takayanagi /
Casini--Huerta knife-edge biconditional that distinguishes
generic-MTC from BH-entropy-yielding boundary conditions
(\cite{ReshetikhinTuraev1991,Witten1989,
NayakSimonSternFreedmanDasSarma2008,FreedmanLarsenWang2002,Kitaev2003,
HaydenPreskill2007,WalkerWang2012,PastawskiYoshidaHarlowPreskill2015}).

\subsection{One MTC, four faces}

The synthesis claim NEW to D4---which the substrate-identity claim
of this flagship inherits---is the observation that the substrate
carries a \emph{single} MTC: the SU(2)$_k$ Chern--Simons MTC at the
horizon. The single MTC underwrites four ostensibly distinct
substrate computations:

\begin{itemize}
\item \textbf{Face 1 (Drinfeld centre).} The Drinfeld-centre
  construction $\mathcal{Z}(\mathrm{Vec}_{\mathbb{Z}_N})$ supplies the
  Standard-Model anomaly classification anchor of Bundle~D2 \S 3.4.
  The cross-bridge here---which Bundle~D4 \S 4 makes explicit---is
  that the centre $\mathbb{Z}_N$ from D3's gauge-erasure mechanism
  (\S\ref{sec:gauge-erasure}) populates the Drinfeld category in the
  toric-code anchor.
\item \textbf{Face 2 (WRT partition function).} The
  Witten--Reshetikhin--Turaev~\cite{ReshetikhinTuraev1991,Witten1989}
  partition function on $S^{2}\times S^{1}$ evaluates to
  $Z(S^{2}\times S^{1}) = \mathrm{rk}\,\mathcal{C}$ (the rank of the
  MTC), and the surgery-formula evaluation on $S^{3}$ delivers a
  closed-form expression for the substrate's BH partition-function
  count of Bundle~D3 \S 7. The trefoil-knot invariant
  $\langle\mathrm{trefoil}\rangle = -\sqrt{2}$ is a load-bearing
  closed-form check.
\item \textbf{Face 3 (Hayden--Preskill QEC).} The horizon-MTC supports
  a Hayden--Preskill quantum error-correcting code structure
  whose code distance and scrambling time satisfy a biconditional
  agreement with the anyonic-fusion-admissibility predicate
  (\texttt{QECHolographyBridge.HPCode.code\_distance\_scaling\_matches\_anyonic\_fusion\_iff\_fusion\_in\_admissible\_class}).
  The substrate's MTC is the same MTC that supplies the
  Bekenstein--Hawking entropy of D3 \S 7 (Face 2) and the boundary-condition
  Prop record \texttt{H\_HorizonBoundaryCondition}.
\item \textbf{Face 4 (RT/Casini--Huerta knife-edge).} The
  Ryu--Takayanagi / Casini--Huerta entanglement-entropy biconditional
  distinguishes generic-MTC boundary conditions from
  BH-entropy-yielding conditions through a knife-edge predicate.
  The same MTC carries both the entropy count of D3 \S 7 and the
  entanglement-entropy structure of the boundary CFT.
\end{itemize}

The unifying claim is that the substrate carries one MTC, not four,
and the four faces are equivalent descriptions of a single
algebraic object on the substrate side. This is the D4 synthesis
claim NEW to D4, which the F flagship lifts here as a substrate-identity
contribution.

\subsection{First-formalisation claims}

D4 ships five first-formalisation claims, each hedged with ``to our
knowledge'' against searches of the Coq, Agda, Isabelle, HOL4, and
Mizar libraries:

\begin{enumerate}
\item $U_q(\mathfrak{sl}_2)$ and $U_q(\mathfrak{sl}_3)$ Hopf-algebra
  axioms with affine and rank-$2$ extensions.
\item Ising MTC braided fusion category with all six hexagon equations
  and ribbon equations, plus the closed-form trefoil invariant
  $\langle\mathrm{trefoil}\rangle = -\sqrt{2}$ and the Hopf-link
  $S$-matrix match.
\item WRT TQFT partition function on $S^{2}\times S^{1}$ and the
  Lickorish surgery formula on three-manifolds (the figure-eight knot
  invariant remains in flight as a tracked-hypothesis closure).
\item Hayden--Preskill QEC code-distance biconditional on horizon-MTC
  boundary conditions, with the Fibonacci-anyon bound as a quantitative
  benchmark
  (\texttt{QECHolographyBridge.fibonacci\_HPCode\_codeDistance\_lt\_log\_two}).
\item Doublon-SWAP gate as the substrate's first symmetry-protected
  topological quantum gate (Berry phase $\pi$, with adiabatic-drag
  premise carried as a tracked-hypothesis Prop record). To our
  knowledge there is no prior published topological quantum gate of
  SPT type in proof-assistant form.
\end{enumerate}

\subsection{Inheritance and axiom budget}

D4 inherits Bundle~D2's single true axiom \texttt{gapped\_interface\_axiom}
(the 4+1D conjecture sitting at the intersection of mechanisms (ii)
and (iii) of D2). \emph{No new axioms are introduced in D4.} The
\texttt{H\_HorizonBoundaryCondition} structure is a tracked-hypothesis
Prop record, not an axiom; D4 \S 7 consumes it through the QEC-MTC
cross-bridge to Bundle~D3 \S 7. The doublon-SWAP gate's adiabatic-drag
premise will be encoded similarly as a tracked-hypothesis Prop record
in the in-flight \texttt{DoublonGate.lean} module.

\subsection{Cross-bridges}

D4 cross-bridges back to three sibling bundles in the substrate-identity
network:

\begin{itemize}
\item \textbf{D4 \S 4 $\leftrightarrow$ D2 \S 3.4 (Drinfeld centre).}
  The toric-code MTC anchor of D4 \S 4 is the same algebraic structure
  that Bundle~D2 \S 3.4 uses as the Standard-Model anomaly classification
  Drinfeld-centre construction. Cross-bridge verified character-for-character
  in the Phase~7c.8 review pass.
\item \textbf{D4 \S 7 $\leftrightarrow$ D3 \S 7, \S 9 (HP-QEC + horizon
  boundary condition).} The Hayden--Preskill code structure consumes
  D3's \texttt{H\_HorizonBoundaryCondition} Prop record; the
  scrambling-time bound and code-distance biconditional are computed
  on the same MTC structure.
\item \textbf{D4 \S 8 $\leftrightarrow$ D3 \S 7 (RT/CH knife-edge).}
  The Casini--Huerta knife-edge biconditional distinguishes
  generic-MTC boundary conditions from BH-entropy-yielding conditions
  on the same horizon MTC that anchors D3's BH entropy.
\end{itemize}

The cross-bridge network closes self-consistently: all three sibling
bundles cite the same Lean theorem names (verified
character-for-character in the Phase~7c.8 review pass), and the
\texttt{H\_HorizonBoundaryCondition} Prop record is defined exactly
once (in D3's \texttt{BHEntropyMicroscopic.lean}) and consumed by
D4 through the cross-bridge.

%% =====================================================================
%% §8 — Dark sector under substrate constraints (D5)
%% =====================================================================
\section{Dark sector under substrate constraints}
\label{sec:dark-sector}

%% TODO: 10pp lift from Bundle D5 (full content).
%% Key claims:
%%   - Six Phase-5x DM mechanism classification.
%%   - Phase 6m Track A: causal-set DE NO-GO + 3 publishable structural caveats.
%%   - Phase 6m Track B: entropic-gravity unanimous NO-GO closure.
%%   - Phase 6m Track C: Jacobson-thermo-GR with M3 f(R) Exp+ArcTanh strongest CLEARED-R5.
%%   - Sakharov 4-condition criterion: first systematic Lambda_J vs Lambda_HK comparison.
%%   - Barrow-bound prescription dependence (embedded structural-caveat subsection 10.5).

\textit{Substantive content awaiting all-12-siblings-GREEN gate (D5 already closed).}

%% =====================================================================
%% §9 — Methodology and the verification stack (I1 + I2)
%% =====================================================================
\section{Methodology and the verification stack}
\label{sec:methodology}

Bundle~I1 (verification methodology, with worked cases) and Bundle~I2
(verified statistics + Lean tensor-categories software paper) describe
the substrate program's verification machinery. Both bundles are
infrastructure-grade content rather than direct substrate physics:
they explain how the substrate program reaches the predictive-register
conclusions that the deep papers (D1--D5) report.

\subsection{The three-layer verification architecture}

Every load-bearing claim in the SK--EFT Hawking program is verified
through three independent layers whose disagreement constitutes
evidence that something is wrong with the substrate analysis:

\begin{enumerate}
\item \textbf{Layer 1: Python numerics.} Quantitative predictions
  ($T_H$ values, dispersion-correction percentages, gain thresholds,
  cosmological-constant ratios, regime-partition mass scales) are
  computed in canonical Python modules (\texttt{src/core/formulas.py},
  \texttt{src/core/constants.py}, \texttt{src/wkb/spectrum.py},
  \texttt{src/adw/gap\_equation.py}). Pipeline Invariant \#1
  designates \texttt{formulas.py} as canonical---the only place
  physics formulas live.
\item \textbf{Layer 2: Lean 4 formal proofs.} Every formula in
  \texttt{formulas.py} carries a Lean theorem reference; every Lean
  theorem has zero \texttt{sorry}. The Lean library currently ships
  $\sim \totalsubstantivetheorems$ substantive theorems across
  \totalmodules{} modules with $\axioms$ axiom and $\sorries$
  \texttt{sorry}.
\item \textbf{Layer 3: Aristotle automated theorem prover.}
  Hard-to-discharge sub-lemmas surfaced during a Lean development
  cycle can be submitted to the Aristotle automated theorem prover
  (\cite{Roehm2026Strategy}, paper15 / Bundle~I1), which closes them
  on the Aristotle-pinned toolchain ($\sim 322$ Aristotle-proved
  theorems contribute to the project total). The Aristotle layer is
  optional but reduces the manual proof effort substantially on
  intricate algebraic obstructions.
\end{enumerate}

The architecture's point is robustness, not redundancy: when the
three layers disagree (as they have in tracked occurrences), the
disagreement \emph{is} the evidence. Bundle~I1 documents three worked
cases of layer-disagreement-as-evidence: (i)~the FirstOrderKMS
Aristotle counterexample, where Aristotle returned a counter-witness
that disproved a tracked Lean conjecture in flight, redirecting the
proof architecture; (ii)~the gap-solution-bounded conjecture disproof,
where a Lean attempt to discharge a Python numerical claim revealed
the claim was inconsistent with a separate Lean theorem;
(iii)~the chirality-wall axiom decomposition, where the original
single-axiom approach was decomposed into three pillars after a
Lean-development cycle showed the single-axiom version was not
discharge-able.

\subsection{The 14-stage Wave Execution Pipeline}

The program ships content through a 14-stage pipeline documented in
\texttt{docs/WAVE\_EXECUTION\_PIPELINE.md} and described in detail
in Bundle~I1 \S 6:

\begin{enumerate}\setlength{\itemsep}{0pt}
\item Constants and refs (canonical: \texttt{constants.py}).
\item Formulas (canonical: \texttt{formulas.py}).
\item Lean (canonical: \texttt{lean/SKEFTHawking/}).
\item Aristotle (optional layer; closes hard sub-lemmas).
\item Lean build verification (\texttt{lake build}).
\item Tests (\texttt{pytest}; canonical: \texttt{tests/}).
\item Validate (\texttt{validate.py}; 22+ checks).
\item Visualizations (canonical: \texttt{visualizations.py}).
\item Figure review (\texttt{physics-qa:figure-reviewer}).
\item Paper draft (\texttt{papers/<bundle>/paper\_draft.tex}).
\item Notebooks (companion narrative content where applicable).
\item Documentation sync (counts, inventory, roadmap, README).
\item Adversarial review (\texttt{adversarial-reviewer} agent;
  fresh-context Stage-13 pass).
\item Quality-improvement register (\texttt{docs/QI\_REGISTER.md};
  process-level lessons banked for future waves).
\end{enumerate}

The pipeline is enforced by Pipeline Invariants 1--14 documented in
\texttt{CLAUDE.md} and Bundle~I1 \S 7. The most load-bearing
invariants for the substrate-identity claim of this flagship are:
Invariant \#1 (formulas.py canonical), Invariant \#5 (every Lean
theorem has a proof), Invariant \#10 (no maxHeartbeats in proof
bodies; decomposition over compute), Invariant \#14 (sentence-level
\texttt{bundle\_destination} tags carry over from source papers).

\subsection{The pre-emptive-strengthening discipline}

Bundle~I1 \S 7 describes a pre-emptive-strengthening discipline that
applies five-question audit to each new theorem statement before
writing it: (i)~drop redundant conjuncts; (ii)~quantitative connection
to numerical content; (iii)~cross-module bridge integrity;
(iv)~trivial-discharge check; (v)~defining-the-conclusion check. The
discipline reduces the post-wave strengthening cost from $\sim 12$
retroactive theorems per wave (the 6c.3 baseline) to $\sim 5$
($58\%$ reduction in 6b.1, 2026-04-27). Empirically, the discipline
catches the more obvious failure modes (existential-absorption,
biconditional-tautology in conditional violators) but \emph{misses}
subtler patterns (identity-function wrappers,
within-own-$\pm 2\sigma$-band tautologies, pairwise-distinctness on
inductive constructors, definitional-unfolding-as-physics).

\subsection{The fresh-context adversarial-reviewer pattern}

Bundle~I1 \S 9 documents the adversarial-reviewer pattern: every
bundle's Stage-13 review is performed by a fresh-context reviewer
agent that has not participated in the bundle's drafting. The fresh
context is what surfaces issues that an authoring-context reviewer
would miss (the authoring context already accepts the bundle's
implicit assumptions). The reviewer's tool budget is hard-capped
($\le 25$ reads, $\le 40$ tool calls) to enforce a finite-context
walk; partial-coverage caveats are documented in the verdict file
when the cap is hit.

The 13-bundle architecture is gated by reviewer-triple closure (Stage
9 figure-reviewer + Stage 10 claims-reviewer + Stage 13
adversarial-reviewer; all three GREEN for a bundle to be marked
closed). At F flagship close-out, all 12 sibling bundles plus this
flagship will have passed reviewer-triple closure independently.

\subsection{Algebraic-Lorentzian-geometry first-formalisations}

Bundle~I1 \S 10 is the primary first-formalisation claim for the
substrate's algebraic-Lorentzian-geometry backbone (Phase 6f / 6g).
Bundle~D3 \S 22.5 is the supplement that consumes these on the
substrate side: D3 cites the I1 first-formalisation primary while
D3 extends the construction to the singularity, area, and no-hair
theorems on the substrate-specific Lorentzian-metric typeclass.

To our knowledge, this is the first machine-checked package on
indefinite-signature metrics in any proof assistant for: the
Lorentzian-metric typeclass, the algebraic first Bianchi identity,
the differential second Bianchi identity, the causal-structure axiom
hierarchy, and the Penrose / Hawking--Penrose / area / Kerr-family
no-hair theorems. The Mathlib upstream cycle target (Bundle~I2) will
contribute the no-hair theorem and the algebraic first Bianchi
identity; the rest is Mathlib-PR-quality but not yet upstream.

\subsection{Bundle I2: VerifiedJackknife and Lean tensor-categories}

Bundle~I2~\cite{Roehm2026Strategy} ships the verified-statistics
software paper content. The two anchors are: (i)~VerifiedJackknife,
a Lean implementation of bias-corrected jackknife statistical estimators
that closes the substrate's RHMC-Monte-Carlo error-bar discipline at
the Lean-formalisation scope; (ii)~the Lean tensor-categories
upstream-Mathlib cycle, which contributes a substantial fraction of
the substrate's category-theoretic infrastructure (Drinfeld centre,
ribbon equations, hexagon equations, MTC braiding) to the public
Mathlib mathematical library. The Mathlib upstream cycle is on the
Phase 8 outlook roadmap; if the cycle is delayed beyond the Phase 7
drafting window, I2 ships as software-only with a JOSS-update target
on the upstream-cycle close.

%% =====================================================================
%% §10 — Substrate predictive register: the boundary partition
%% =====================================================================
\section{Substrate predictive register: the boundary partition}
\label{sec:substrate-register}

This section consolidates the substrate's complete predictive register
across the boundary partition introduced in \S\ref{sec:strategic-frame}.
Each entry traces to a specific Lean theorem or experimental anchor
in a sibling bundle. The register is the substrate program's
substantive output: \emph{what the substrate predicts at quantitative
precision} versus \emph{what the substrate fails to predict at
quantitative precision}.

\subsection{Cleared register}

The substrate produces quantitative predictions agreeing with
observation within the published experimental band at the following
locations:

\begin{itemize}
\item \textbf{Decision Gate~E.1 (linearised EFE).} The Vergeles
  positivity / critical-limit / Adler-reduction triple discharges
  cleanly at the Sakharov-identity calibration point on the
  substrate's tracked-hypothesis Prop record bundle (D3
  \S\ref{sec:emergent-gravity} with Bundle~D3 \S 2).
\item \textbf{Decision Gate~E.2 (heat-kernel calibration).} The
  Sakharov coefficient $G_N^{\mathrm{Sak}} = 12\pi/(N_f\,\Lambda_{\mathrm{UV}}^{2})$
  is verified character-for-character against Vassilevich
  Eq.~4.37~\cite{Vassilevich2003}. The same coefficient propagates
  across six observables in D3 without an additional dial.
\item \textbf{Decision Gate~E.3 (diff invariance through $a_4$).}
  The path-(b) order-by-order invariance failure
  \texttt{NonlinearDiffInvariance.pathB\_residual\_a4\_dirac\_eq\_zero}
  closes on the Dirac fermion bundle; perturbations away from the
  Dirac basis introduce a residual linear in the perturbation
  parameter (D3 \S 19).
\item \textbf{Decision Gate~E.4 (multi-channel PPN).} The
  precession-ratio biconditional
  \texttt{NonlinearEFE.precessionRatio\_eq\_one\_iff\_alpha\_unity}
  closes the multi-channel PPN comparison against
  Will~2014~\cite{Will2014} (D3 \S 20).
\item \textbf{Polariton stimulated-Hawking gain.} $G > 0.5$ at
  $\omega/\kappa = \ln 3 / (2\pi) \approx 0.175$ with
  $N_{\mathrm{probe}} = 100$ probe photons per mode for $5\sigma$
  detection. Bundle~D1 (\S\ref{sec:analog-hawking}) plus Bundle~E1
  experimental letter for the Paris--LKB platform.
\item \textbf{Graphene Dirac-fluid analog Hawking.} $T_H \approx 2.4\,$K,
  $\delta_{\mathrm{disp}} \approx -2.8\%$, Wiedemann--Franz violation
  $L/L_0 > 200$. Bundle~D1 plus Bundle~E2 experimental letter for the
  Dean--Kim--Lucas platform.
\item \textbf{Phase 6m Track-C $f(R)$ Exp + ArcTanh.} The strongest
  CLEARED-R5 of any track in the Phase 6m three-track closure
  (Plaza--Kraiselburd:
  $\Delta\!\mathrm{AIC} \simeq \Delta\!\mathrm{BIC} \gtrsim 20$,
  Jeffreys ``very strong''; D5 \S\ref{sec:dark-sector}).
\item \textbf{Einstein--Cartan torsion at the cosmological background.}
  $|T_{\mathrm{EC}}| \approx 2.05\times 10^{-77}\,\mathrm{GeV}$, which
  is $\sim 46$ orders below the
  Kosteleck\'y--Russell--Tasson~\cite{KosteleckyRussellTasson2008}
  experimental bound (D3 \S 22).
\item \textbf{Stelle higher-curvature ceilings.} The substrate's
  predicted contributions to the Stelle $(\alpha,\beta,\gamma)$
  coefficients sit safely below the LIGO + pulsar-timing-array +
  binary-BH-inspiral observational ceilings (D3 \S 18).
\end{itemize}

\subsection{Falsified register}

The substrate produces predictions in tension with observation at
$\geq 5\sigma$ or Bayesian thresholds exceeding decisive at the
following locations:

\begin{itemize}
\item \textbf{GW170817 vestigial-second-sound graviton.} The
  vestigial-mode-as-graviton hypothesis is falsified at $\sim 7\times
  10^{14}$ over the Volovik vestigial natural range (D3
  \S\ref{sec:emergent-gravity} with Bundle~L1 PRL splash). The
  falsification is universal across the natural-range parameter band
  $\chi_{\mathrm{vest}}\in[0.1,10]$ via Abbott et al.\
  $|\Delta c/c| \le 3\times 10^{-15}$~\cite{Abbott2017GW170817}.
\item \textbf{Wen--ADW factor-$6000$ deficit.} The
  perturbative-Wen extension to gravity under-predicts $G_N$ by a
  closed-form Lean ratio $G_{\mathrm{Wen}}/G_c^{\mathrm{ADW}} \approx
  1/6000$ (D3 \S 4 with~\cite{Wen2003}). The deficit closes the
  emergent-non-Abelian-gauge-symmetry pathway to gravity at
  quantitative magnitude.
\item \textbf{Fracton diffeomorphism-invariance no-go.} The
  Pretko--Radzihovsky fracton hydrodynamics
  channel~\cite{PretkoRadzihovsky2018} fails to support full
  diffeomorphism invariance at order $a_4$; the obstruction shows up
  exactly at the order at which Stelle higher-curvature
  $(\alpha,\beta,\gamma)$ live (D3 \S 4).
\item \textbf{Abelian-MTC Bekenstein--Hawking falsifier.} The
  substrate's MTC-derived $-3/2 \log A$ correction to the
  Bekenstein--Hawking entropy~\cite{KaulMajumdar2000} is falsified
  on abelian MTC choices (the falsifier pinpoints SU(2)$_k$ as the
  unique MTC consistent with both the entropy count and the QEC
  structure of D4). D3 \S 7 with D4 \S\ref{sec:topological-qc}.
\item \textbf{Sen $-45/2$ disagreement with Kaul--Majumdar.} The
  Sen 4D Schwarzschild logarithmic-correction coefficient $-45/2$
  disagrees with the Kaul--Majumdar SU(2)$_k$ $-3/2$ at quantitative
  precision; the substrate's choice of MTC is structurally selected
  by the falsifier (D3 \S 7).
\item \textbf{Phase 6m Track-B entropic-gravity dark-energy NO-GO
  closure.} The first complete-mechanism-family unanimous NO-GO
  closure in Phase~6m: $8/8$ candidates closed with
  $|\!\log\mathcal{B}| \geq 5$ (Jeffreys-decisive); $7/8$ candidates
  $r_d$-independent (D5 \S\ref{sec:dark-sector}).
\item \textbf{Hu--Sawicki $f(R)$ chameleon Solar-System NO-GO.} The
  Hu--Sawicki variant of Track C is closed via the
  Solar-System chameleon constraint at $b \approx 0.21$
  ($> 100\times$ the chameleon bound; D5 \S\ref{sec:dark-sector}).
\item \textbf{Pure Lovelock $\alpha_2$ NO-GO.} The Pure Lovelock
  candidate is closed at the $1\sigma$-box edge of Quintom-B
  ($|w_0+1| < 0.05$, $|\tilde\alpha_2|_{\max}\le 0.15$;
  D5 \S\ref{sec:dark-sector}).
\item \textbf{EWBG doubly-forbidden.} The substrate's electroweak
  baryogenesis pathway is doubly forbidden by the chirality-wall
  + sphaleron-rate combination (D3 \S 13 with~\cite{KlinkhamerManton1984,
  ArnoldMcLerran1987}).
\end{itemize}

\subsection{Structurally-blocked register}

The substrate's mechanism is inapplicable for a load-bearing structural
reason at the following locations:

\begin{itemize}
\item \textbf{Gibbs--Duhem obstruction theorem.} Any single-scalar
  self-tuning emergent-vacuum framework with Gibbs--Duhem equilibrium
  locks $w_{\mathrm{vac}} = -1$ (D5 \S\ref{sec:dark-sector}). The
  obstruction is realisation-independent across the four tested KV
  q-theory constructions (Bundle~D5; Phase 5y).
\item \textbf{Four-factor orthogonality principle.} The dark-energy
  obstruction follows from four orthogonal factors:
  $\text{GD} \cap c_s^{2}\!\geq\!0 \cap \text{natural }T_c \cap
  \text{MICROSCOPE}$. Vestigial gravity fails through factor \#1
  (Gibbs--Duhem); other Volovik-family mechanisms fail through other
  factor combinations (D5).
\item \textbf{Causal-set DE GD-inapplicability across all four
  sprinkling prescriptions.} Phase 6m Track-A (D5
  \S\ref{sec:dark-sector}). Three publishable structural caveats
  survive independent of DESI: GD-inapplicability across
  prescriptions; Barrow-bound prescription dependence; BDG
  $\sigma_\Lambda = \alpha_{\mathrm{BDG}}/\sqrt{V}$ first-principles
  decomposition.
\item \textbf{Higher-form-symmetry erasure.} Only $U(1)$ survives the
  fluid layer; non-Abelian one-form symmetries are structurally
  erased by hydrodynamic coarse-graining (D3 \S 3 with
  \S\ref{sec:gauge-erasure}).
\end{itemize}

\subsection{Reproduced (not solved) register}

The substrate produces predictions matching observation at the
structural level but failing to solve the corresponding fine-tuning
problem at the following locations:

\begin{itemize}
\item \textbf{Cosmological-constant problem.} The substrate's
  Sakharov coefficient predicts an effective $\Lambda$ that
  overshoots the observed value by
  $\Lambda_{\mathrm{emerg}}/\Lambda_{\mathrm{obs}} \gtrsim 10^{100}$
  (rigorous Lean lower bound; heuristically $\sim 10^{122}$).
  D3 \S 21. The substrate inherits the $a_0$ heat-kernel
  coefficient's universality across emergent-gravity proposals at
  one-loop; emergent gravity does not shelter from the
  cosmological-constant problem~\cite{Weinberg1989}.
\end{itemize}

\subsection{Open under tracked hypothesis}

The substrate's predictive register includes obligations that are
explicitly open under tracked-hypothesis Lean Props, with closure paths
documented in each sibling bundle's Outlook section:

\begin{itemize}\setlength{\itemsep}{0pt}
\item Explicit one-loop calculation of $\alpha_{\mathrm{ADW}}$
  discharging the three Vergeles propositions (D3 \S 2 + \S 5).
\item RHMC convergence at $L \geq 8$ in the vestigial-metric Monte
  Carlo (D3 \S 10).
\item Higgs-mass back-reaction calculation
  closing the BHL bilocal mechanism (D3 \S 11).
\item PMNS structure on the Majorana rung (D3 \S 12).
\item Figure-eight knot invariant
  in the Ising MTC braided fusion category (D4 \S 4).
\item Full WRT surgery formula for arbitrary closed orientable
  3-manifolds (D4 \S 5; partial coverage on $S^3$ via Lickorish
  surgery).
\item Doublon-SWAP gate adiabatic-drag premise
  (D4 \S 6, in flight in \texttt{DoublonGate.lean}).
\item Mathlib upstream cycle for the no-hair theorem and the
  algebraic first Bianchi identity (I2; Phase 8 outlook).
\item Phase 6m Track-C M8 KSS path-(c) universal-horizon
  thermodynamics (D5 \S\ref{sec:dark-sector}).
\end{itemize}

The open-under-tracked-hypothesis register is itself a substrate
prediction: each entry is a quantitatively well-defined obligation
whose closure is on the Phase 8+ roadmap. The substrate program is
not done; it is at a coherent intermediate state where every closure
target is identified.

%% =====================================================================
%% §11 — Open problems and Phase 8 outlook
%% =====================================================================
\section{Open problems and Phase~8 outlook}
\label{sec:open-problems}

\subsection{Phase 7 close-out summary}

At the time of this flagship's submission, the substrate program has
shipped 13 bundles to reviewer-triple-closed Stage-13 GREEN status.
The 12 sibling bundles ship across five tiers: Tier~0 (this flagship),
Tier~1 (D1--D5; five deep papers), Tier~2 (L1--L3; three PRL splashes
extracted from D2 and D3), Tier~3 (I1--I2; two infrastructure papers),
and Tier~4 (E1--E2; two experimental letters). The architecture is
the project's strategy frame~\cite{Roehm2026Strategy}.

Total program output at flagship close-out:
$\sim \totalsubstantivetheorems$ substantive Lean theorems across
\totalmodules{} modules with $\axioms$ axiom and $\sorries$ \texttt{sorry}.
The single axiom is \texttt{gapped\_interface\_axiom}, the 4+1D
gapped-interface conjecture from Bundle~D2's chirality-wall pillar 2.
All other load-bearing claims are theorems.

\subsection{Submission roll-out}

The submission sequence per~\cite{Roehm2026Strategy} \S 3 begins with
the Tier-2 PRL splashes (L1, L2, L3) as arXiv vouchers ahead of the
Tier-1 deep-paper submissions. The Tier-1 deep papers (D1, D2, D3,
D4, D5) follow the splashes, with D3 (the heaviest at $\sim 50$pp)
shipped after L1 and L3 since L1 and L3 extract content from D3.
The Tier-3 infrastructure papers (I1, I2) ship in parallel with the
Tier-1 deep papers; I2 may delay if the Mathlib upstream cycle
extends. The Tier-4 experimental letters (E1, E2) ship with the D1
deep paper as companion content. This flagship review (F) ships last,
citing all 12 siblings in their Phase-7-final form.

The submission cadence is intentionally tied to the architecture's
internal cross-bridges: each downstream submission cites upstream
submissions in their final-published form rather than as in-preparation
self-cites. The full submission roll-out is on the Phase 8 roadmap.

\subsection{Phase 6 deferred targets}

The substrate program closed Phase 6m (the three-track dark-energy
deep-research closure) at 2026-04-30. Phase 6 has produced
sub-phases 6a through 6m to date. Several Phase 6 deferred targets
remain on the Phase 8 outlook roadmap:

\begin{itemize}
\item \textbf{Phase 6n+ anticipated.} Future Phase 6X waves may ship
  during Phase 7 / Phase 8 execution. The
  \texttt{LATE\_PHASE6\_ABSORPTION\_PROTOCOL.md} (a frozen 7-stage
  protocol with branches D.1/D.2/D.3) provides the canonical
  workflow for absorbing late-Phase-6 outputs into already-drafted
  bundles without redrafting.
\item \textbf{Strong-CP / Z16 / theta-vacuum Phase 6l/6m extension.}
  The substrate is silent on the strong-CP problem at the structural
  level: $\bar\theta$ is uniquely orthogonal to the substrate's
  $\mathbb{Z}_{16}$ anomaly classification (Phase 6l Wave 1 verdict).
  This isolates strong-CP as the unique BSM-hierarchy substrate
  $\mathbb{Z}_{16}$ does not fix.
\item \textbf{Phase 6f / 6g substrate roster algebraic-Lorentzian-geometry
  first-formalisations.} Bundle~I1 \S 10 ships these as the primary
  first-formalisation claim; D3 \S 22.5 is the supplement on the
  substrate side.
\end{itemize}

\subsection{Mathlib upstream cycle}

Bundle~I2's Mathlib upstream cycle is the single largest piece of
infrastructure work on the Phase 8 roadmap. The contributions targeted
for upstream are: the no-hair theorem on the Lorentzian-metric
typeclass; the algebraic first Bianchi identity on bundle Riemann
tensors; the ribbon equations and hexagon equations of the Ising MTC
braided fusion category; the WRT TQFT partition-function evaluation
on $S^{2}\times S^{1}$; the Drinfeld-centre construction on the
toric-code anchor. The cycle target is to reduce the substrate
program's reliance on project-internal duplications of standard
mathematical material; success of the cycle would let future
substrate-program work import these directly from Mathlib rather than
ship them in-bundle.

If the cycle is delayed beyond the Phase 7 drafting window, Bundle~I2
ships as software-only with a JOSS-update target on the upstream-cycle
close.

\subsection{Outlook and architectural roadmap}

Looking forward beyond Phase 7, the substrate program's architectural
roadmap rests on four open questions whose answers are on Phase 8+
horizons:

\begin{enumerate}
\item \textbf{Is the substrate's predictive register complete at the
  Lean-formalisation scope?} The current state ($\axioms$ axiom,
  $\sorries$ \texttt{sorry}) is consistent with completeness, but the
  open-under-tracked-hypothesis register
  (\S\ref{sec:substrate-register}) lists $\sim 10$ residual closure
  targets. Each is bounded; the question is whether closure of all of
  them produces qualitatively new content or only confirms the
  current substrate-identity claim.
\item \textbf{Does the experimental program (Bundle~E1 polariton
  Paris--LKB and Bundle~E2 graphene Dean--Kim--Lucas) produce a
  positive observation?} The substrate's predictions are quantitative;
  the Tier-4 experimental letters give the device-parameter audit.
  Phase 8 collaboration with the Paris--LKB and Dean--Kim--Lucas
  groups is on the roadmap.
\item \textbf{Is the chirality-wall axiom dischargeable?} The single
  remaining axiom \texttt{gapped\_interface\_axiom} is a mathematical
  conjecture (existence of a particular gapped-interface
  construction). Its discharge is a Mathlib-upstream-cycle target
  but may also depend on independent mathematical-topology research.
\item \textbf{What is the substrate's response to a positive DESI~DR3
  $f(R)$ measurement?} The strongest CLEARED-R5 candidates of
  Phase 6m Track-C ($f(R)$ Exp + ArcTanh) are the likely DESI~DR3
  test target; a positive measurement would clear them at the
  predictive-register cleared category, while a negative measurement
  would close the only remaining alternative-mechanism family at
  Phase 6m.
\end{enumerate}

The substrate program is at a coherent intermediate state. Its
architectural commitment is to deliver the substrate's predictive
register at quantitative precision across the four boundary-partition
categories without rescoping in the face of negative results.

%% =====================================================================
%% §12 — Conclusions
%% =====================================================================
\section{Conclusions}
\label{sec:conclusions}

\subsection{The substrate is one object}

The synthesis claim of this flagship is structurally absent from any
individual sibling bundle. Each sibling reports its own thread; this
flagship reports the substrate as a single object across all twelve.
The cross-thread identities listed in \S\ref{sec:strategic-frame} are
the F-unique synthesis content:

\begin{itemize}
\item The $N_f$ in
  $G_N^{\mathrm{Sak}} = 12\pi/(N_f\,\Lambda_{\mathrm{UV}}^{2})$
  (D3 \S\ref{sec:emergent-gravity}) and the $N_f = 16$ that uniquely
  fixes the Standard-Model anomaly classification via
  $\mathrm{Spin}^{\mathbb{Z}_4}$ 5-bordism (D2
  \S\ref{sec:chirality-wall-modular} with L2) are the same $N_f$.
\item The SK--EFT axioms, KMS condition, and fluctuation-dissipation
  relation that close the BEC analog Hawking spectrum (D1
  \S\ref{sec:analog-hawking}) re-close on polariton (E1) and
  graphene Dirac-fluid (E2) platforms with $\sim 92\%$ Lean-theorem
  reuse: the substrate is platform-agnostic at the formal level.
\item The MTC carrying the Kaul--Majumdar $-3/2 \log A$ correction to
  Bekenstein--Hawking entropy (D3 \S\ref{sec:emergent-gravity}) is
  the same MTC supporting Hayden--Preskill recovery on the horizon
  (D4 \S\ref{sec:topological-qc}) and the Drinfeld-centre cross-bridge
  to the Standard-Model anomaly classification (D2 with D4): one MTC,
  four faces.
\item The Sakharov-induced-gravity criterion that distinguishes
  substrate classes within Phase~6m's Track C dark-energy taxonomy (D5
  \S\ref{sec:dark-sector}) is the same criterion that classifies
  substrate-vs-non-substrate analog systems
  ($^3$He-A vs.\ acoustic-BEC); to our knowledge this is the first
  systematic $\Lambda_J$-vs-$\Lambda_{HK}$ comparison on a common
  substrate in the literature.
\end{itemize}

The substrate is one object; its twelve faces are thirteen bundles'
worth of predictive register.

\subsection{Predictive register at the boundary partition}

The substrate's predictive register lives at the boundary partition
of phenomenology---each prediction lands cleared, falsified,
structurally blocked, or reproduced. This is what distinguishes
substrate physics from speculative model-building: the substrate
fails predictively at definite locations, and the substrate succeeds
predictively at definite locations, and both kinds of result are
load-bearing.

The cleared register includes the four cleared decision gates of D3,
the polariton stimulated-Hawking-gain prediction of D1+E1, and the
strongest Phase~6m Track-C $f(R)$ candidates of D5. The falsified
register includes the GW170817 vestigial-graviton falsification of
D3+L1 at $\sim 7\times 10^{14}$, the Wen--ADW factor-$6000$ deficit
of D3, the abelian-MTC entropy falsifier of D3+D4, and the unanimous
$8/8$ Phase~6m Track-B closure of D5. The structurally-blocked
register includes the Gibbs--Duhem obstruction theorem, the
fracton diff-invariance no-go, and the four-factor orthogonality
principle. The reproduced register includes the cosmological-constant
problem at $\Lambda_{\mathrm{emerg}}/\Lambda_{\mathrm{obs}} \gtrsim
10^{100}$ (rigorous Lean lower bound) or heuristically $\sim
10^{122}$.

\S\ref{sec:substrate-register} ships the consolidated table.

\subsection{The 13-bundle architecture as a coherent submission package}

Per the project's strategy frame~\cite{Roehm2026Strategy}, the 13
bundles ship as a coherent submission package: the five Tier-1 deep
papers carry the full substantive content; the three Tier-2 PRL
splashes extract headline results from D3 (L1, L3) and D2 (L2); the
two Tier-3 infrastructure papers describe the methodology (I1) and
the verified-statistics software (I2); the two Tier-4 experimental
letters carry per-platform companion content (E1 polariton, E2
graphene); and this Tier-0 flagship review (Bundle~F) integrates the
twelve sibling threads into a single survey with substrate-identity
synthesis claim. NO-GO results are reported as first-class predictive
content rather than absences. All load-bearing claims ship with Lean
4 modules at zero \texttt{sorry} and one true axiom
(\texttt{gapped\_interface\_axiom}; the 4+1D gapped-interface
conjecture). The architecture is intentionally append-friendly: new
Tier-2 PRL splashes can be added if a future Phase 6X wave produces
a 4-page-PRL-shaped headline.

The substrate is one object; the 13-bundle architecture is its
coherent description across SM-plus-GR phenomenology.

%% BUNDLE_APPEND_INSERT_HERE — bundle_append.py inserts new sections at
%% this marker. To insert at a specific section, edit body manually after
%% the script runs.

\begin{thebibliography}{99}

%% F bibliography lifts citations from all 12 sibling bundles. The
%% inline-thebibliography style (rather than \bibliography{<file>.bib})
%% avoids the empty-.bib failure mode that surfaced for D3 + D4 in
%% Phase 7c.6 (see qi-bundle_skeleton_inline_bibliography in
%% docs/QI_REGISTER.md).

%% TODO: populate ~150-200 \bibitem{} entries by lifting the union of
%% sibling-bundle bibliographies once all 12 siblings reach Stage-13
%% GREEN. CITATION_REGISTRY membership is the single source of truth;
%% bibitem text must match registry entry character-for-character per
%% qi-bibitem_registry_drift.

\bibitem{Roehm2026Strategy}
J.~G.~Roehm,
``SK--EFT Hawking 13-bundle publication architecture,''
project paper-strategy frame, in preparation (2026).

%% Sibling-bundle self-cites (in-prep until each ships):
\bibitem{Roehm2026D1}
J.~G.~Roehm,
``Analog Hawking radiation across three platforms (Bundle D1),''
in preparation (2026).

\bibitem{Roehm2026D2}
J.~G.~Roehm,
``Anomaly constraints on Standard Model particle content (Bundle D2),''
in preparation (2026).

\bibitem{Roehm2026D3}
J.~G.~Roehm,
``Emergent gravity through black-hole thermodynamics (Bundle D3),''
in preparation (2026).

\bibitem{Roehm2026D4}
J.~G.~Roehm,
``Topological quantum-computation foundations: Drinfeld centre,
Hayden--Preskill, and the longest verified mathematical chain
(Bundle D4),''
in preparation (2026).

\bibitem{Roehm2026D5}
J.~G.~Roehm,
``Dark sector under substrate constraints (Bundle D5),''
in preparation (2026).

\bibitem{Roehm2026E1}
J.~G.~Roehm,
``Polariton stimulated-Hawking gain at modest probe-photon counts
(Bundle E1 letter),''
in preparation (2026).

\bibitem{Roehm2026E2}
J.~G.~Roehm,
``Graphene Dirac-fluid analog Hawking radiation: device-parameter
audit (Bundle E2 letter),''
in preparation (2026).

%% TODO: lift all external citations from sibling bundles (D1-D5,
%% L1-L3, I1-I2, E1-E2). Approximately 150-200 unique \bibitem{}
%% entries expected after the 12-of-13 gate clears.

\end{thebibliography}

\end{document}
